% META: {"id": "TNSI-ARCH-EVAL-B", "chapitre": "TNSI-ARCHITECTURES-MATERIELLES-SY", "type_objet": "evaluation", "capacites": ["T-ARCH-03", "T-ARCH-04A", "T-ARCH-04B"], "mode_creation": "ex_nihilo", "sources_inspiration": [], "duree_min": 45, "status": "generated"}
\section*{Évaluation B --- Routage et cryptographie}
\textit{Durée : 45 minutes. Ordinateur avec interpréteur Python autorisé.}

\baremeIndicatif{Ex. 1 : 10 pts (routage) ; Ex. 2 : 10 pts (chiffrement, HTTPS)}

\bigskip

\textbf{Exercice 1} --- \textit{Routage} \hfill (10 points)

\begin{python}
reseau = {
    "A": {"B": 1, "D": 8},
    "B": {"A": 1, "C": 1},
    "C": {"B": 1, "D": 1},
    "D": {"A": 8, "C": 1},
}
\end{python}
\begin{enumerate}
  \item (5 pts) En utilisant \lstinline{bfs_sauts} du cours, quelle route RIP choisirait-il de \texttt{A} vers \texttt{D}, et en combien de sauts ?
  \item (5 pts) En utilisant \lstinline{plus_court_cout} du cours, quel est le coût de la route la moins coûteuse de \texttt{A} vers \texttt{D} ? Coïncide-t-elle avec la route RIP ?
\end{enumerate}

% BEGIN-VERIFY
% def bfs_sauts(graphe_couts, depart, arrivee):
%     visites = {depart: 0}
%     file = [depart]
%     while file:
%         courant = file.pop(0)
%         if courant == arrivee:
%             return visites[courant]
%         for voisin in graphe_couts[courant]:
%             if voisin not in visites:
%                 visites[voisin] = visites[courant] + 1
%                 file.append(voisin)
%     return None
% def plus_court_cout(graphe_couts, depart, arrivee):
%     couts = {depart: 0}
%     a_traiter = [depart]
%     while a_traiter:
%         courant = min(a_traiter, key=lambda s: couts[s])
%         a_traiter.remove(courant)
%         if courant == arrivee:
%             break
%         for voisin, poids in graphe_couts[courant].items():
%             nouveau_cout = couts[courant] + poids
%             if voisin not in couts or nouveau_cout < couts[voisin]:
%                 couts[voisin] = nouveau_cout
%                 if voisin not in a_traiter:
%                     a_traiter.append(voisin)
%     return couts[arrivee]
% reseau = {
%     "A": {"B": 1, "D": 8},
%     "B": {"A": 1, "C": 1},
%     "C": {"B": 1, "D": 1},
%     "D": {"A": 8, "C": 1},
% }
% assert bfs_sauts(reseau, "A", "D") == 1
% assert plus_court_cout(reseau, "A", "D") == 3
% END-VERIFY

\bigskip

\textbf{Exercice 2} --- \textit{Chiffrement et HTTPS} \hfill (10 points)

\begin{enumerate}
  \item (4 pts) Quelle est la différence essentielle entre chiffrement symétrique et asymétrique ?
  \item (6 pts) Expliquer pourquoi une connexion HTTPS combine les deux, plutôt que d'utiliser uniquement le chiffrement asymétrique pour tout l'échange.
\end{enumerate}
